\chapter{William P. Thurston --- Fields Medal 1982}\label{ch1982williamthurston:william-p.-thurston-fields-medal-1982-a-comprehensive-report}

\begin{center}\rule{0.5\linewidth}{0.5pt}\end{center}

\textbf{Prepared:} July 2026 \textbf{Subject:} William P. Thurston (b.
October 30, 1946, Washington, D.C., USA; d.~August 21, 2012)
\textbf{Position:} Professor, University of California, Berkeley;
Cornell University; Princeton University \textbf{Doctoral Advisor:}
Morris Hirsch (University of California, Berkeley) \textbf{Fields Medal
Citation:} \emph{``For his significant contributions to the study of
3-manifolds and the geometrization conjecture.''}

\begin{center}\rule{0.5\linewidth}{0.5pt}\end{center}

\section{1. Introduction}\label{ch1982williamthurston:introduction}

William P. Thurston was awarded the Fields Medal in 1982 for his
substantial work in low-dimensional topology, particularly his
\textbf{geometrization conjecture} for 3-manifolds. This conjecture,
which he formulated in the late 1970s and proved in many cases himself,
proposed that every compact 3-manifold can be decomposed into pieces
that each admit one of eight canonical geometric structures. This
transformed the study of 3-manifolds from a largely topological subject
into one closely connected with geometry, analysis, and dynamics.

Thurston's contributions span a notable range of mathematics: -
The geometrization conjecture, which classifies all 3-manifolds by their
geometric structures and was eventually proved in full by Perelman in
2003 using Ricci flow. - The classification of surface diffeomorphisms
and the study of pseudo-Anosov mappings, which revolutionized the study
of surface dynamics and Teichmüller theory. - The theory of hyperbolic
3-manifolds, showing that ``most'' 3-manifolds are hyperbolic and
providing the first systematic construction of hyperbolic structures. -
The development of the orbifold theorem, proved by Thurston and later by
Boileau--Leeb--Porti. - The Thurston norm on the homology of 3-manifolds
and its relationship to foliations. - Fundamental contributions to the
theory of foliations, including the Thurston stability theorem. - The
study of combinatorial geometry and the development of circle packing
techniques. - The definition and study of the stretch factor (Thurston's
algorithm) for surface automorphisms. - Deep contributions to the theory
of Kleinian groups and Mostow rigidity.

Thurston's geometric vision and his ability to see deep connections
between disparate areas of mathematics made him one of the most
influential mathematicians of the late 20th century. His work has had a
lasting impact on topology, geometry, dynamics, and group theory.

This report provides a comprehensive overview of Thurston's life, work,
and lasting impact on mathematics, organized by the major themes of his
research.

\begin{center}\rule{0.5\linewidth}{0.5pt}\end{center}

\section{2. Biographical Background}\label{ch1982williamthurston:biographical-background}

\subsection{2.1 Early Life and
Education}\label{ch1982williamthurston:early-life-and-education}

William Paul Thurston was born in 1946 in Washington, D.C. He showed
mathematical talent from an early age, building topological models as a
teenager. He entered New College in Sarasota, Florida, where he studied
mathematics. He then moved to the University of California, Berkeley for
graduate studies.

\subsection{2.2 Doctoral Work}\label{ch1982williamthurston:doctoral-work}

At Berkeley, Thurston studied under Morris Hirsch, working on foliations
of 3-manifolds. His PhD thesis (1972) on foliations already contained
groundbreaking ideas, including the Thurston stability theorem for
foliations.

\subsection{2.3 Academic Career}\label{ch1982williamthurston:academic-career}

\begin{itemize}
\tightlist
\item
 \textbf{1973:} PhD, University of California, Berkeley.
\item
 \textbf{1973--1974:} Member, Institute for Advanced Study, Princeton.
\item
 \textbf{1974--1991:} Professor, Princeton University --- where he
 developed the geometrization program.
\item
 \textbf{1991--1996:} Professor, University of California, Berkeley.
\item
 \textbf{1996--2003:} Professor, University of California, Davis.
\item
 \textbf{2003--2012:} Professor, Cornell University.
\item
 \textbf{1982:} Fields Medal awarded at the ICM in Warsaw.
\item
 \textbf{2012:} Died at age 65 from complications of melanoma.
\end{itemize}

\begin{center}\rule{0.5\linewidth}{0.5pt}\end{center}

\section{3. The Geometrization
Conjecture}\label{ch1982williamthurston:the-geometrization-conjecture}

\subsection{3.1 Background}\label{ch1982williamthurston:background}

Before Thurston, the classification of 3-manifolds was a major open
problem, one of the important in topology. For surfaces
(2-manifolds), the uniformization theorem provides a complete
classification: every closed surface admits a Riemannian metric of
constant curvature \(+1\), \(0\), or \(-1\). This geometric
classification underlies the entire theory of Riemann surfaces. For
3-manifolds, no such unified classification existed, despite the work of
Poincaré, Haken, Waldhausen, and others.

\subsection{3.2 Thurston's Conjecture}\label{ch1982williamthurston:thurstons-conjecture}

Thurston conjectured that every compact 3-manifold can be decomposed
along incompressible tori and spheres into pieces, each of which admits
one of eight geometric structures. These eight geometries are: 1.
\textbf{Spherical} \(S^3\): constant positive curvature, corresponding
to elliptic geometry. 2. \textbf{Euclidean} \(\mathbb{E}^3\): constant
zero curvature, corresponding to flat geometry. 3. \textbf{Hyperbolic}
\(\mathbb{H}^3\): constant negative curvature, the richest and most
common case. 4. \(S^2 \times \mathbb{R}\): the product of a sphere and a
line. 5. \(\mathbb{H}^2 \times \mathbb{R}\): the product of the
hyperbolic plane and a line. 6. \(\widetilde{\text{SL}_2\mathbb{R}}\):
the universal cover of the Lie group \(\text{SL}_2\mathbb{R}\). 7.
\textbf{Nil} geometry: the Heisenberg group, a nilpotent Lie group. 8.
\textbf{Sol} geometry: the solvable Lie group, the least symmetric of
the eight.

\subsection{3.3 The Eight Geometries}\label{ch1982williamthurston:the-eight-geometries}

Each of the eight geometries is a model geometry \((X, G)\) where \(X\)
is a simply connected Riemannian manifold and \(G\) is a Lie group
acting transitively on \(X\) with compact stabilizers. A 3-manifold
admits a geometric structure modeled on \((X, G)\) if it is
diffeomorphic to \(X/\Gamma\) for a discrete subgroup
\(\Gamma \subset G\) acting freely.

\subsection{3.4 Significance}\label{ch1982williamthurston:significance}

The geometrization conjecture is the 3-dimensional analogue of the
uniformization theorem for surfaces. It provides a complete
classification of 3-manifolds and reduces many topological questions to
geometric ones. For example: - The Poincaré conjecture (every simply
connected closed 3-manifold is \(S^3\)) follows from the spherical case.
- The Borel conjecture about rigidity of aspherical manifolds reduces to
the hyperbolic case. - Many topological invariants of 3-manifolds can be
computed from their geometric decomposition.

\subsection{3.5 Proof by Perelman}\label{ch1982williamthurston:proof-by-perelman}

The full geometrization conjecture was proved by Grigori Perelman in
2003 using Richard Hamilton's Ricci flow. Perelman showed that Ricci
flow with surgery converges to a collection of geometric pieces,
confirming Thurston's vision. Perelman's proof earned him a Fields Medal
(which he declined) and a Clay Millennium Prize. The proof substantially
relies on Thurston's geometric insight and his decomposition framework.

\begin{center}\rule{0.5\linewidth}{0.5pt}\end{center}

\section{4. Hyperbolic 3-Manifolds}\label{ch1982williamthurston:hyperbolic-3-manifolds}

\subsection{4.1 Mostow Rigidity}\label{ch1982williamthurston:mostow-rigidity}

Mostow's rigidity theorem states that a closed hyperbolic 3-manifold is
determined up to isometry by its fundamental group. This means that the
hyperbolic structure is unique when it exists --- topological invariants
become geometric invariants. This rigidity, which fails in higher
dimensions, makes hyperbolic 3-manifolds particularly tractable.

\subsection{4.2 Thurston's Hyperbolization
Theorem}\label{ch1982williamthurston:thurstons-hyperbolization-theorem}

Thurston proved that a Haken 3-manifold (one containing an
incompressible surface) that is atoroidal admits a hyperbolic structure.
This was the first major evidence for the geometrization conjecture and
remains one of its deepest partial results. The proof uses: - The theory
of Kleinian groups and their deformation spaces (quasiconformal
deformations). - The measurable Riemann mapping theorem and the
Ahlfors--Bers theory. - The analysis of the boundaries of deformation
spaces. - The construction of pleated surfaces and the use of the
Thurston norm.

\subsection{4.3 The Density of Hyperbolic
3-Manifolds}\label{ch1982williamthurston:the-density-of-hyperbolic-3-manifolds}

Thurston showed that ``most'' 3-manifolds are hyperbolic through his
work on Dehn surgery. The \textbf{hyperbolic Dehn surgery theorem}
states that all but finitely many Dehn surgeries on a hyperbolic knot
yield hyperbolic 3-manifolds. Moreover, the set of hyperbolic
3-manifolds is dense in the space of all 3-manifolds in a precise sense.

\subsection{4.4 Volumes of Hyperbolic
3-Manifolds}\label{ch1982williamthurston:volumes-of-hyperbolic-3-manifolds}

Thurston studied the volume function on the space of hyperbolic
3-manifolds, showing that: - The volumes of hyperbolic 3-manifolds form
a well-ordered subset of \(\mathbb{R}\). - There is a smallest volume
hyperbolic 3-manifold (conjectured to be the Weeks manifold). - The
volume is related to the simplicial volume (Gromov norm) and satisfies
the proportionality principle.

\begin{center}\rule{0.5\linewidth}{0.5pt}\end{center}

\section{5. The Thurston Norm}\label{ch1982williamthurston:the-thurston-norm}

\subsection{5.1 Definition}\label{ch1982williamthurston:definition}

The \textbf{Thurston norm} on \(H_2(M, \partial M)\) measures the
minimal complexity of a surface representing a given homology class. For
a compact, orientable 3-manifold \(M\) and an embedded surface \(S\)
with Euler characteristic \(\chi(S)\), define: \[
x(S) = \max\{-\chi(S), 0\}
\] For a homology class \(a \in H_2(M, \partial M)\), the Thurston norm
is: \[
x(a) = \min\{x(S) : S \text{ is an embedded surface representing } a\}
\]

\subsection{5.2 Properties}\label{ch1982williamthurston:properties}

The Thurston norm extends to a seminorm on the real homology
\(H_2(M, \partial M; \mathbb{R})\). Its unit ball is a convex polytope
whose vertices are integer homology classes. The norm detects whether a
homology class can be represented by a surface of a given Euler
characteristic.

\subsection{5.3 Relationship to
Foliations}\label{ch1982williamthurston:relationship-to-foliations}

The Thurston norm is intimately connected to the theory of taut
foliations. Thurston proved that the unit ball of the Thurston norm is a
polytope whose faces correspond to taut foliations on the 3-manifold.
Specifically, a homology class \(a\) is represented by a taut foliation
if and only if it lies on the boundary of the unit ball and its
supporting face is rational.

\begin{center}\rule{0.5\linewidth}{0.5pt}\end{center}

\section{6. Foliations}\label{ch1982williamthurston:foliations}

\subsection{6.1 The Thurston Stability
Theorem}\label{ch1982williamthurston:the-thurston-stability-theorem}

Thurston's stability theorem is a fundamental result in foliation
theory. It states that a foliation of a compact manifold with no
holonomy --- meaning that nearby leaves are not ``twisted'' around each
other --- must have a compact leaf. More precisely, if \(\mathcal{F}\)
is a foliation of a compact manifold \(M\) such that the holonomy group
of every leaf is trivial, then there exists a compact leaf.

\subsection{6.2 Codimension-One
Foliations}\label{ch1982williamthurston:codimension-one-foliations}

Thurston made deep contributions to the study of codimension-one
foliations on 3-manifolds: - He showed that every closed 3-manifold with
\(b_1(M) > 0\) admits a codimension-one foliation. - He developed the
relationship between foliations and the Thurston norm, showing that taut
foliations correspond to faces of the Thurston norm ball. - He proved
that the existence of a taut foliation implies that the 3-manifold is
covered by \(\mathbb{R}^3\). - The construction of foliations from
branched surfaces and the study of their transverse measures. -
Applications of foliations to the study of mapping tori and surface
bundles over the circle.

\begin{center}\rule{0.5\linewidth}{0.5pt}\end{center}

\section{7. Orbifolds}\label{ch1982williamthurston:orbifolds}

\subsection{7.1 Definition}\label{ch1982williamthurston:definition-1}

An \textbf{orbifold} is a generalization of a manifold that allows for
cone points (points with isotropy) and reflector edges. Orbifolds arise
naturally in the study of 3-manifolds as quotients of manifolds by
finite group actions. The geometrization conjecture extends naturally to
orbifolds, which are essential for understanding branched covers and
Seifert fibered spaces.

\subsection{7.2 The Orbifold Theorem}\label{ch1982williamthurston:the-orbifold-theorem}

Thurston proved the orbifold theorem: every compact, orientable
3-orbifold with a geometric decomposition admits a geometric structure.
This theorem was later rigorized and generalized by Boileau, Leeb, and
Porti (2005). The orbifold theorem is essential for: - Understanding
branched covers of 3-manifolds. - The classification of finite group
actions on 3-manifolds. - The proof of the Smith conjecture and other
classical problems.

\begin{center}\rule{0.5\linewidth}{0.5pt}\end{center}

\section{8. Surface Automorphisms and Pseudo-Anosov
Maps}\label{ch1982williamthurston:surface-automorphisms-and-pseudo-anosov-maps}

\subsection{8.1 Classification of Surface
Diffeomorphisms}\label{ch1982williamthurston:classification-of-surface-diffeomorphisms}

Thurston's classification of surface diffeomorphisms (often called the
\textbf{Nielsen--Thurston classification}) is a fundamental result
in low-dimensional topology. It divides all elements of the mapping
class group of a compact surface into three mutually exclusive types: -
\textbf{Periodic:} \(f^n = \text{id}\) for some \(n > 0\), corresponding
to finite-order automorphisms of a Riemann surface. -
\textbf{Reducible:} \(f\) preserves a multicurve (a disjoint union of
essential simple closed curves), reducing the complexity of the surface.
- \textbf{Pseudo-Anosov:} \(f\) expands and contracts along a pair of
transverse measured foliations, the most common and most interesting
type.

\subsection{8.2 Pseudo-Anosov Maps}\label{ch1982williamthurston:pseudo-anosov-maps}

A pseudo-Anosov map \(f: S \to S\) preserves a pair of transverse
measured foliations \((\mathcal{F}^u, \mathcal{F}^s)\) and stretches one
by a factor \(\lambda > 1\) (the \textbf{dilatation} or stretch factor)
while contracting the other by \(1/\lambda\). These maps are the
building blocks of surface dynamics and exhibit: - Exponential growth of
the number of periodic points. - Mixing dynamics with a Markov
partition. - Minimality of the invariant foliations. - Relationship to
Teichmüller theory, where pseudo-Anosov maps correspond to geodesics in
Teichmüller space.

\subsection{8.3 The Stretch Factor}\label{ch1982williamthurston:the-stretch-factor}

The stretch factor \(\lambda > 1\) of a pseudo-Anosov map is a key
invariant. Thurston showed that \(\lambda\) is an algebraic integer
whose degree is bounded by the topology of the surface. The set of
possible stretch factors for a given surface is well-ordered, with the
smallest being the \textbf{minimal dilatation}.

\begin{center}\rule{0.5\linewidth}{0.5pt}\end{center}

\section{9. Kleinian Groups}\label{ch1982williamthurston:kleinian-groups}

\subsection{9.1 Definition}\label{ch1982williamthurston:definition-2}

A \textbf{Kleinian group} is a discrete subgroup
\(\Gamma \subset \text{PSL}_2\mathbb{C}\) acting on the Riemann sphere
\(\widehat{\mathbb{C}}\) by Möbius transformations. The quotient
\(\mathbb{H}^3/\Gamma\) is a hyperbolic 3-manifold (or orbifold).
Kleinian groups are thus the fundamental groups of hyperbolic
3-manifolds, and their study connects 3-manifold topology with complex
analysis, dynamics, and group theory.

\subsection{9.2 Thurston's Contributions}\label{ch1982williamthurston:thurstons-contributions}

Thurston made fundamental contributions to the theory of Kleinian
groups: - \textbf{Deformation theory:} Thurston developed a complete
theory of deformations of Kleinian groups, parameterizing hyperbolic
structures on 3-manifolds by the space of quasiconformal deformations. -
\textbf{Density of cusps:} Thurston showed that cusped hyperbolic
3-manifolds are dense in the deformation space of hyperbolic structures.
- \textbf{Geometric limits:} The study of limits of Kleinian groups
under algebraic and geometric convergence, leading to the theory of
strong convergence. - \textbf{The ending lamination theorem:} Thurston
formulated the ending lamination conjecture (proved by
Brock--Canary--Minsky), which classifies Kleinian groups by their
topological type and ending laminations. - \textbf{The Bers--Thurston
theory:} Connections between Kleinian groups, Teichmüller theory, and
the boundary of deformation spaces.

\begin{center}\rule{0.5\linewidth}{0.5pt}\end{center}

\section{10. Circle Packing}\label{ch1982williamthurston:circle-packing}

\subsection{10.1 Definition}\label{ch1982williamthurston:definition-3}

\textbf{Circle packing} is the arrangement of circles in the plane (or
on a surface) such that each pair of circles is either disjoint or
externally tangent. A circle packing is determined by its
\textbf{nerve}, a graph whose vertices correspond to circles and edges
to tangencies. The \textbf{circle packing theorem}
(Koebe--Andreev--Thurston) states that every planar graph can be
realized as the nerve of a circle packing.

\subsection{10.2 The Thurston--Rodin--Sullivan
Theorem}\label{ch1982williamthurston:the-thurstonrodinsullivan-theorem}

Thurston conjectured that circle packings could approximate conformal
maps, providing a discrete analogue of the Riemann mapping theorem. This
was proved by Rodin and Sullivan (1987): given a simply connected region
\(\Omega \subset \mathbb{C}\) and a regular hexagonal circle packing of
\(\Omega\), the packing converges to the Riemann mapping as the circle
radii go to zero. This opened the field of \textbf{discrete complex
analysis} and has applications to: - Numerical conformal mapping. - The
study of circle packings on surfaces and 3-manifolds. - Connections to
hyperbolic geometry and the study of 3-manifolds.

\begin{center}\rule{0.5\linewidth}{0.5pt}\end{center}

\section{11. The Work on 3-Manifold
Invariants}\label{ch1982williamthurston:the-work-on-3-manifold-invariants}

\subsection{11.1 The Thurston--Bennequin
Number}\label{ch1982williamthurston:the-thurstonbennequin-number}

In contact geometry and knot theory, the Thurston--Bennequin number
\(tb(K)\) is an invariant of Legendrian knots in contact 3-manifolds. It
measures the twisting of the contact framing relative to the Seifert
framing along the knot. The Thurston--Bennequin number satisfies the
Thurston--Bennequin inequality and plays a central role in the
classification of Legendrian and transverse knots.

\subsection{11.2 Simplicial Volume}\label{ch1982williamthurston:simplicial-volume}

Thurston defined the \textbf{simplicial volume} (also called the Gromov
norm) \(\|M\|\) of a closed manifold \(M\). It is defined as the infimum
of \(\sum |a_i|\) over all singular chains representing the fundamental
class of \(M\). For hyperbolic 3-manifolds, the simplicial volume is
proportional to the hyperbolic volume: \[
\|M\| = \frac{\text{Vol}(M)}{v_3}
\] where \(v_3\) is the volume of the ideal regular tetrahedron in
\(\mathbb{H}^3\). This proportionality theorem is one of the
foundational results in the study of hyperbolic manifolds.

\begin{center}\rule{0.5\linewidth}{0.5pt}\end{center}

\section{12. The Work on the Mapping Class
Group}\label{ch1982williamthurston:the-work-on-the-mapping-class-group}

\subsection{12.1 The Mapping Class Group}\label{ch1982williamthurston:the-mapping-class-group}

The mapping class group \(\text{Mod}(S)\) of a surface \(S\) is the
group of isotopy classes of orientation-preserving diffeomorphisms of
\(S\). One of the central problems in low-dimensional topology is
understanding the structure of this group.

\subsection{12.2 The Nielsen--Thurston
Classification}\label{ch1982williamthurston:the-nielsenthurston-classification}

Thurston's classification of surface automorphisms provides a complete
understanding of the mapping class group of a surface. The
Nielsen--Thurston classification divides mapping classes into three
types: periodic, reducible, and pseudo-Anosov. This classification is
the analogue for surfaces of the classification of elements of
\(\text{SL}_2\mathbb{Z}\) into elliptic, parabolic, and hyperbolic
types.

\subsection{12.3 The Thurston Measure}\label{ch1982williamthurston:the-thurston-measure}

Thurston introduced a natural measure \(\mu_{\text{Th}}\) on the space
\(\mathcal{MF}\) of measured foliations on a surface. The Thurston
measure is invariant under the action of the mapping class group and is
used to compute intersection numbers and to study the dynamics of
pseudo-Anosov maps. The total mass of the Thurston measure on the
projectivized space \(\mathcal{PMF}\) is finite and related to the
Weil--Petersson volume of moduli space.

\begin{center}\rule{0.5\linewidth}{0.5pt}\end{center}

\section{13. The Work on Combinatorial
Geometry}\label{ch1982williamthurston:the-work-on-combinatorial-geometry}

\subsection{13.1 Thick-Thin
Decomposition}\label{ch1982williamthurston:thick-thin-decomposition}

A fundamental concept in hyperbolic geometry is the thick-thin
decomposition of a hyperbolic manifold \(M\). The \textbf{thin part}
\(M_{\text{thin}}\) consists of points where the injectivity radius is
less than a fixed constant \(\epsilon\) (the Margulis constant). For
hyperbolic 3-manifolds, the thin part consists of: - \textbf{Cusps:}
Tubular neighborhoods of short geodesics that are actually tori cusps
(for manifolds with torus boundary). - \textbf{Tubes:} Solid tori around
short geodesics of length less than \(\epsilon\). The \textbf{thick
part} \(M_{\text{thick}}\) has bounded geometry and is compact when
\(M\) has finite volume.

\subsection{13.2 The Jørgensen--Thurston
Theorem}\label{ch1982williamthurston:the-juxf8rgensenthurston-theorem}

Jørgensen and Thurston studied the space of hyperbolic 3-manifolds with
finite volume and showed that: - The volume function
\(\text{Vol}: \mathcal{M}_3 \to \mathbb{R}_{\geq 0}\) is continuous in
the topology of geometric convergence. - There are only finitely many
hyperbolic 3-manifolds with volume less than any given bound. - The set
of volumes is well-ordered, and the minimal volume hyperbolic
3-manifolds occur at the bottom of this set.

\subsection{13.3 The Volume Spectrum}\label{ch1982williamthurston:the-volume-spectrum}

The study of volumes of hyperbolic 3-manifolds is a major area of
research. The smallest known orientable hyperbolic 3-manifold is the
Weeks manifold (volume \(\approx 0.9427\)). It is conjectured to be the
unique minimal volume hyperbolic 3-manifold.

\begin{center}\rule{0.5\linewidth}{0.5pt}\end{center}

\section{14. The Work on Nielsen
Theory}\label{ch1982williamthurston:the-work-on-nielsen-theory}

\subsection{14.1 The Nielsen Realization
Problem}\label{ch1982williamthurston:the-nielsen-realization-problem}

The Nielsen realization problem asks whether every finite subgroup of
the mapping class group \(\text{Mod}(S)\) can be realized as a group of
automorphisms of some Riemann surface structure on \(S\). Thurston
solved this problem in the affirmative, using his work on the geometry
and dynamics of surface automorphisms.

\subsection{14.2 Proof Strategy}\label{ch1982williamthurston:proof-strategy}

Thurston's proof uses: - The Nielsen--Thurston classification of mapping
classes. - The action of the mapping class group on Teichmüller space
and its Thurston compactification. - Fixed point theory and the Brouwer
fixed point theorem. - The fact that finite subgroups have a global
fixed point in the Thurston compactification, which corresponds to a
Riemann surface structure where the group acts by isometries.

\subsection{14.3 Applications}\label{ch1982williamthurston:applications}

The Nielsen realization theorem has important applications to: - The
study of moduli spaces of Riemann surfaces and their orbifold structure.
- The structure of the mapping class group and its subgroups. - The
theory of Teichmüller spaces and the geometry of the mapping class group
action.

\begin{center}\rule{0.5\linewidth}{0.5pt}\end{center}

\section{15. The Work on the Smith
Conjecture}\label{ch1982williamthurston:the-work-on-the-smith-conjecture}

\subsection{15.1 Statement}\label{ch1982williamthurston:statement}

The Smith conjecture states that a periodic diffeomorphism of the
3-sphere \(S^3\) has a fixed point set that is an unknotted circle. This
seemingly simple statement was one of the major open problems in
3-manifold topology for many decades.

\subsection{15.2 Thurston's Contribution}\label{ch1982williamthurston:thurstons-contribution}

Thurston contributed to the proof of the Smith conjecture as part of a
collaborative effort in the late 1970s. The proof used: - Thurston's
work on hyperbolic geometry and the orbifold theorem. - The theory of
Kleinian groups and their deformation spaces. - The classification of
finite group actions on 3-manifolds. The solution of the Smith
conjecture was a landmark result that demonstrated the power of
Thurston's geometric methods.

\subsection{15.3 Related Results}\label{ch1982williamthurston:related-results}

The methods developed for the Smith conjecture were subsequently applied
to prove other classical results, including the characterization of
finite group actions on \(S^3\) and the classification of spherical
space forms.

\begin{center}\rule{0.5\linewidth}{0.5pt}\end{center}

\section{16. The Work on the Marden
Conjecture}\label{ch1982williamthurston:the-work-on-the-marden-conjecture}

\subsection{16.1 Statement}\label{ch1982williamthurston:statement-1}

The Marden conjecture (now the Marden--Thurston theorem, proved by Agol
and Calegari--Gabai) states that a hyperbolic 3-manifold with finitely
generated fundamental group is \textbf{tame} --- that is, it is
homeomorphic to the interior of a compact 3-manifold with boundary.

\subsection{16.2 Thurston's
Contribution}\label{ch1982williamthurston:thurstons-contribution-1}

Thurston proved the Marden conjecture for Haken manifolds as part of his
hyperbolization theorem. More specifically, he showed that any
hyperbolic 3-manifold that is geometrically finite and contains an
incompressible surface is tame. The general case was later proved by: -
Calegari and Gabai (2006) for the case of freely indecomposable
fundamental group. - Agol (2004) for the general case, completing the
proof.

\subsection{16.3 Significance}\label{ch1982williamthurston:significance-1}

The tameness theorem is fundamental for the structure theory of
hyperbolic 3-manifolds. It implies that every hyperbolic 3-manifold with
finitely generated fundamental group has a compact core (the Scott core
theorem with a geometric refinement) and can be studied using the
methods of geometric topology.

\begin{center}\rule{0.5\linewidth}{0.5pt}\end{center}

\section{17. The Work on the Volume
Conjecture}\label{ch1982williamthurston:the-work-on-the-volume-conjecture}

\subsection{17.1 Statement}\label{ch1982williamthurston:statement-2}

The \textbf{volume conjecture} (formulated by Kashaev and later
Murakami--Murakami) relates the hyperbolic volume of a knot complement
to the asymptotic behavior of the colored Jones polynomial.
Specifically, for a hyperbolic knot \(K\), the conjecture states: \[
\lim_{n \to \infty} \frac{2\pi \log |J_n(K, e^{2\pi i/n})|}{n} = \text{Vol}(S^3 \setminus K)
\] where \(J_n(K, q)\) is the \(n\)-colored Jones polynomial.

\subsection{17.2 Thurston's
Contribution}\label{ch1982williamthurston:thurstons-contribution-2}

Thurston's work on hyperbolic knot complements provided the geometric
foundation for the volume conjecture: - He showed that the complement of
a hyperbolic knot admits a complete hyperbolic structure of finite
volume. - He developed the theory of ideal triangulations of hyperbolic
knot complements, which allows for the computation of volumes. - His
work on the relationship between geometry and quantum invariants
inspired the formulation of the volume conjecture. - The conjecture has
been verified for many knots and is one of the important
connections between geometry and quantum topology.

\begin{center}\rule{0.5\linewidth}{0.5pt}\end{center}

\section{18. The Work on Cobordisms}\label{ch1982williamthurston:the-work-on-cobordisms}

\subsection{18.1 Thurston's Work}\label{ch1982williamthurston:thurstons-work}

Thurston studied the cobordism theory of 3-manifolds and its
relationship to foliations and the Thurston norm. The cobordism group
\(\Omega_3\) classifies 3-manifolds up to cobordism (the existence of a
4-manifold bounding them). Thurston showed that: - The Thurston norm is
a key tool in understanding the cobordism group of 3-manifolds. - Taut
foliations give rise to nontrivial cobordism invariants. - The
relationship between the Thurston norm and the cobordism group leads to
constraints on which 3-manifolds can bound a 4-manifold with a given
signature.

\begin{center}\rule{0.5\linewidth}{0.5pt}\end{center}

\section{19. The Work on the Geometrization of
2-Complexes}\label{ch1982williamthurston:the-work-on-the-geometrization-of-2-complexes}

\subsection{19.1 Curvature of
2-Complexes}\label{ch1982williamthurston:curvature-of-2-complexes}

Thurston studied geometric structures on 2-dimensional complexes and
their relationship to 3-manifold topology. This work contributed to: -
The theory of non-positively curved metric spaces and the work of Gromov
on hyperbolic groups. - The study of CAT(0) geometries and their
applications to geometric group theory. - The understanding of how
geometric structures on 2-complexes relate to the topology of
3-manifolds branched over them.

\begin{center}\rule{0.5\linewidth}{0.5pt}\end{center}

\section{20. The Work on the Measured
Foliations}\label{ch1982williamthurston:the-work-on-the-measured-foliations}

\subsection{20.1 Definition}\label{ch1982williamthurston:definition-4}

A \textbf{measured foliation} \(\mathcal{F}\) on a surface \(S\) is a
foliation (a decomposition of \(S\) into 1-dimensional leaves) equipped
with a transverse measure \(\mu\). The transverse measure assigns a
positive number to arcs transverse to the foliation, invariant under
homotopy of the arc. The space \(\mathcal{MF}\) of measured foliations
on \(S\) up to isotopy and Whitehead equivalence is homeomorphic to
\(\mathbb{R}^{6g-6}\) for a surface of genus \(g \geq 2\).

\subsection{20.2 Train Tracks}\label{ch1982williamthurston:train-tracks}

Thurston introduced \textbf{train tracks} as a combinatorial tool for
studying measured foliations and the mapping class group. A train track
is a 1-complex with a smooth structure that models the local structure
of a measured foliation. Train tracks provide: - A finite-dimensional
approximation of the space of measured foliations. - A computational
framework for studying the action of the mapping class group. - A method
for computing stretch factors of pseudo-Anosov maps. Train tracks are
now a standard tool in low-dimensional topology, hyperbolic geometry,
and Teichmüller theory.

\begin{center}\rule{0.5\linewidth}{0.5pt}\end{center}

\section{21. The Work on the Earle--Thurston
Theorem}\label{ch1982williamthurston:the-work-on-the-earlethurston-theorem}

\subsection{21.1 Statement}\label{ch1982williamthurston:statement-3}

The \textbf{Earle--Thurston theorem} (proved with Clifford Earle) states
that the space of conformal structures on a surface of genus
\(g \geq 2\) is homeomorphic to Teichmüller space. This result provides
a topological framework for studying Riemann surfaces and their moduli.

\subsection{21.2 Significance}\label{ch1982williamthurston:significance-2}

The Earle--Thurston theorem is important because: - It shows that
Teichmüller space is topologically a cell, confirming the
contractibility of the space of conformal structures. - It provides a
natural topology for the space of Riemann surfaces that is compatible
with the deformation theory of Kleinian groups. - It is used in the
proof of the Nielsen realization theorem and other results about the
mapping class group.

\begin{center}\rule{0.5\linewidth}{0.5pt}\end{center}

\section{22. The Work on the Thurston
Compactification}\label{ch1982williamthurston:the-work-on-the-thurston-compactification}

\subsection{22.1 Definition}\label{ch1982williamthurston:definition-5}

The \textbf{Thurston compactification} \(\overline{\mathcal{T}}(S)\) of
Teichmüller space \(\mathcal{T}(S)\) adds a boundary
\(\partial\mathcal{T}(S)\) consisting of \textbf{projective measured
foliations} \(\mathcal{PMF}\). The compactification is homeomorphic to a
closed ball of dimension \(6g-6\). The boundary
\(\partial\mathcal{T}(S)\) is homeomorphic to a sphere of dimension
\(6g-7\).

\subsection{22.2 Properties}\label{ch1982williamthurston:properties-1}

The Thurston compactification has the following crucial properties: -
The mapping class group \(\text{Mod}(S)\) acts continuously on the
compactification. - Periodic mapping classes fix interior points of
\(\mathcal{T}(S)\). - Pseudo-Anosov mapping classes have fixed points on
the boundary \(\partial\mathcal{T}(S)\). - Reducible mapping classes fix
points in the interior of boundary strata. - The compactification is
compatible with the Weil--Petersson geometry of Teichmüller space.

\subsection{22.3 Applications}\label{ch1982williamthurston:applications-1}

The Thurston compactification is essential for: - Understanding the
dynamics of the mapping class group action on Teichmüller space. - The
proof of the Nielsen--Thurston classification. - The study of the
boundary of moduli space and its compactification. - The relationship
between Teichmüller theory and the geometry of hyperbolic 3-manifolds.

\begin{center}\rule{0.5\linewidth}{0.5pt}\end{center}

\section{23. The Work on the Markov
Numbers}\label{ch1982williamthurston:the-work-on-the-markov-numbers}

\subsection{23.1 Markov Triples}\label{ch1982williamthurston:markov-triples}

Markov triples are integer solutions to the Markov diophantine equation:
\[
x^2 + y^2 + z^2 = 3xyz
\] Markov numbers are the numbers appearing in Markov triples, and they
form one of the most fascinating sequences in number theory.

\subsection{23.2 Thurston's Connection}\label{ch1982williamthurston:thurstons-connection}

Thurston discovered a deep connection between Markov numbers and
hyperbolic geometry. He showed that Markov numbers correspond to the
cusps of hyperbolic structures on the punctured torus. Specifically: -
Each Markov triple corresponds to a hyperbolic structure on the
complement of a figure-eight knot. - The Markov numbers are related to
the lengths of closed geodesics on the modular torus. - The Markov
spectrum (the set of Markov numbers as a subset of \(\mathbb{R}\)) has
important connections to Diophantine approximation and the Lagrange
spectrum.

\begin{center}\rule{0.5\linewidth}{0.5pt}\end{center}

\section{24. The Work on the Thurston
Algorithm}\label{ch1982williamthurston:the-work-on-the-thurston-algorithm}

\subsection{24.1 Definition}\label{ch1982williamthurston:definition-6}

The Thurston algorithm (also called the \textbf{Thurston train track
algorithm}) computes the stretching factor (dilatation) of a
pseudo-Anosov map by analyzing its action on the space of measured
foliations represented by train tracks.

\subsection{24.2 Algorithm}\label{ch1982williamthurston:algorithm}

The algorithm proceeds as follows: 1. Represent the mapping class by its
action on a train track. 2. Find the induced linear map on the vector
space of weights on the train track. 3. Compute the spectral radius of
this linear map --- this is the stretch factor \(\lambda > 1\). 4. The
corresponding eigenvector gives the invariant measured foliation.

\subsection{24.3 Significance}\label{ch1982williamthurston:significance-3}

The Thurston algorithm is fundamental for: - Computing the dynamics of
surface automorphisms. - Studying the growth of periodic points and
other dynamical invariants. - Understanding the spectrum of possible
dilatations for surfaces of a given genus. - Applications in Teichmüller
theory and the study of the mapping class group.

\begin{center}\rule{0.5\linewidth}{0.5pt}\end{center}

\section{25. Awards and Recognition}\label{ch1982williamthurston:awards-and-recognition}

\textbf{1974:} - Invited Speaker, International Congress of
Mathematicians, Vancouver --- presenting his early work on foliations.

\textbf{1976:} - Alan T. Waterman Award (National Science Foundation)
--- the first mathematician to receive this award, recognizing
outstanding early-career research.

\textbf{1978:} - Invited Plenary Speaker, International Congress of
Mathematicians, Helsinki --- presenting the geometrization program.

\textbf{1982:} - Fields Medal (International Mathematical Union) ---
awarded at the ICM in Warsaw for his significant contributions to the
study of 3-manifolds and the geometrization conjecture.

\textbf{1995:} - Member, National Academy of Sciences (USA).

\textbf{2001:} - Oswald Veblen Prize in Geometry (American Mathematical
Society).

\textbf{2008:} - Leroy P. Steele Prize for Seminal Contribution to
Research (American Mathematical Society) --- for his paper
``Three-dimensional manifolds, Kleinian groups and hyperbolic
geometry.''

\textbf{2012:} - Elected to the American Academy of Arts and Sciences. -
Died at age 65 in Ithaca, New York, from complications of melanoma.

\begin{center}\rule{0.5\linewidth}{0.5pt}\end{center}

\section{26. Selected Major
Publications}\label{ch1982williamthurston:selected-major-publications}

\begin{enumerate}
\def\labelenumi{\arabic{enumi}.}
\item
 \textbf{``The Geometry and Topology of Three-Manifolds.''} Princeton
 lecture notes, 1979. --- \emph{The seminal lecture notes that
 introduced the geometrization conjecture and transformed 3-manifold
 topology.}
\item
 \textbf{``Three-dimensional manifolds, Kleinian groups and hyperbolic
 geometry.''} \emph{Bulletin of the AMS} 6, 357--381, 1982. ---
 \emph{The survey article that laid out the geometrization program.}
\item
 \textbf{``On the geometry and dynamics of diffeomorphisms of
 surfaces.''} \emph{Bulletin of the AMS} 19, 417--431, 1988. ---
 \emph{The classification of surface automorphisms and the introduction
 of pseudo-Anosov maps.}
\item
 \textbf{``The theory of foliations of codimension greater than one.''}
 \emph{Comment. Math. Helv.} 49, 214--231, 1974. --- \emph{The Thurston
 stability theorem for foliations.}
\item
 \textbf{``A norm for the homology of 3-manifolds.''} \emph{Memoirs of
 the AMS} 59, no. 339, 1986. --- \emph{The introduction of the Thurston
 norm and its relationship to foliations.}
\item
 \textbf{``Hyperbolic structures on 3-manifolds I: Deformation of
 acylindrical manifolds.''} \emph{Annals of Mathematics} 124, 203--246,
 1986. --- \emph{The hyperbolization theorem for Haken manifolds.}
\item
 \textbf{``The topology of 3-manifolds.''} \emph{Proceedings of the
 ICM}, Helsinki, 1978. --- \emph{Thurston's ICM address presenting the
 geometrization program.}
\item
 \textbf{``How to see 3-manifolds.''} \emph{Classical and Quantum
 Gravity} 15, 2545, 1998. --- \emph{A survey of 3-manifold topology.}
\item
 \textbf{``On the construction and classification of foliations.''}
 \emph{Proc. Int. Conf. on Prospects in Mathematics}, Kyoto, 1976. ---
 \emph{Contributions to the theory of foliations.}
\item
 \textbf{``Shapes of polyhedra and triangulations of the sphere.''} In
 \emph{The Epstein Birthday Schrift}, 1998. --- \emph{Connections
 between circle packing and 3-manifolds.}
\end{enumerate}

\begin{center}\rule{0.5\linewidth}{0.5pt}\end{center}

\section{27. References and Further
Reading}\label{ch1982williamthurston:references-and-further-reading}

\subsection{Official Sources}\label{ch1982williamthurston:official-sources}

\begin{itemize}
\tightlist
\item
 IMU Fields Medal 1982 Announcement, Warsaw.
\item
 Thurston, W. \emph{The Geometry and Topology of Three-Manifolds}.
 Princeton University lecture notes, 1979. Electronically published by
 MSRI.
\item
 Thurston, W. ``On the geometry and dynamics of diffeomorphisms of
 surfaces.'' \emph{Bull. AMS} 19, 1988.
\end{itemize}

\subsection{Key Technical References}\label{ch1982williamthurston:key-technical-references}

\begin{itemize}
\tightlist
\item
 Thurston, W. \emph{The Geometry and Topology of Three-Manifolds}.
 Princeton lecture notes, 1979.
\item
 Thurston, W. ``Three-dimensional manifolds, Kleinian groups and
 hyperbolic geometry.'' \emph{Bull. AMS} 6, 357--381, 1982.
\item
 Thurston, W. ``Hyperbolic structures on 3-manifolds I: Deformation of
 acylindrical manifolds.'' \emph{Ann. Math.} 124, 203--246, 1986.
\item
 Thurston, W. ``A norm for the homology of 3-manifolds.'' \emph{Memoirs
 AMS} 59, 1986.
\item
 Morgan, J. ``On Thurston's uniformization theorem for
 three-dimensional manifolds.'' In \emph{The Smith Conjecture},
 Academic Press, 1984.
\item
 McMullen, C. ``Riemann surfaces and the geometrization of
 3-manifolds.'' \emph{Bull. AMS} 27, 1992.
\end{itemize}

\subsection{Textbooks}\label{ch1982williamthurston:textbooks}

\begin{itemize}
\tightlist
\item
 Benedetti, R. and Petronio, C. \emph{Lectures on Hyperbolic Geometry}.
 Springer Universitext, 1992.
\item
 Canary, R., Epstein, D., and Marden, A. (eds.) \emph{Fundamentals of
 Hyperbolic Manifolds}. Cambridge University Press, 2006.
\item
 Kapovich, M. \emph{Hyperbolic Manifolds and Discrete Groups}.
 Birkhäuser, 2001.
\item
 Martelli, B. \emph{An Introduction to Geometric Topology}.
 CreateSpace, 2016.
\item
 Ratcliffe, J. \emph{Foundations of Hyperbolic Manifolds}. Springer GTM
 149, 2006.
\item
 Farb, B. and Margalit, D. \emph{A Primer on Mapping Class Groups}.
 Princeton University Press, 2012.
\end{itemize}

\subsection{Biography and Laudations}\label{ch1982williamthurston:biography-and-laudations}

\begin{itemize}
\tightlist
\item
 Gabai, D. ``William P. Thurston.'' \emph{Notices of the AMS}, 60(10),
 2013.
\item
 Sullivan, D. ``The Work of William Thurston.'' ICM 1982 Laudatio.
\item
 Morgan, J. ``The Geometrization Conjecture.'' In \emph{Fields
 Medalists' Laudations}, IMU, 1982.
\item
 Thurston, W. ``William P. Thurston.'' In \emph{The Princeton Companion
 to Mathematics}, 2008.
\end{itemize}

\begin{center}\rule{0.5\linewidth}{0.5pt}\end{center}

\emph{(End of report)}
